% id: 202-15
% book: 베이직쎈 대수 · 12 수학적 귀납법
% section: 유형 125 수학적 귀납법
% figure: none
% answer: 
% answer_source: 
% variant_level: 0 (원본 전사)
% 이 파일은 build.mjs 가 items.json 에서 만든다. 고칠 때는 items.json 을 고친다.
\smash{\raisebox{1.5mm}{\dmkichul{베이직쎈 대수 202쪽 15번}}}
자연수 $n$에 대한 명제 $p(n)$이 모든 짝수에 대하여 성립함을 수학적 귀납법을 이용하여 증명하려면 다음 두 가지를 보이면 된다.\exprbox{\begin{minipage}{0.96\linewidth}\raggedright ㈎ $p(a)$가 참이다.\\[2.5mm] ㈏ $p(k)$가 참이면 $p(k+b)$도 참이다.\end{minipage}}이때 자연수 $a$, $b$에 대하여 $ab$의 값을 구하시오.
